\section{Model and data details}\label{app:details}

\subsection{Forecast systems}
Table~\ref{tab:models} lists the systems. All reported results are out of
sample with respect to model development. IFS, AIFS, and AIFS ENS changed
operational versions on 12 May 2026 \citep{ecmwf2026cycle50r1}; the analysis
treats each named product as one system and does not isolate version effects.

\begin{table}[h]
\caption{Forecast systems. Ensemble systems are evaluated through their
means. ICON-EU appears only in Track B; Helios only in the solar comparison.}
\label{tab:models}
\centering\small
\resizebox{\textwidth}{!}{%
\begin{tabular}{llll}
\toprule
Model & Class & Output & Domain \\
\midrule
ECMWF IFS (reference) & physical & deterministic & global \\
ECMWF ENS & physical & ensemble & global \\
NOAA GFS & physical & deterministic & global \\
DWD ICON Global & physical & deterministic & global \\
DWD ICON-EU & physical & deterministic & regional \\
ECMWF AIFS & regression AI & deterministic & global \\
ECMWF AIFS ENS & proper-score AI & 51-member ensemble & global \\
Microsoft Aurora & regression AI & deterministic & global \\
Jua EPT-2 Reasoning & regression AI & deterministic & global \\
Jua EPT-2e & regression AI & ensemble & global \\
Jua EPT-2 HRRR & generative AI & ensemble & regional \\
Jua EPT-2.1 Europa & generative AI & ensemble & regional \\
Jua EPT-2.1 Helios & solar AI & deterministic & regional \\
\bottomrule
\end{tabular}}
\end{table}

\subsection{Station panel and tracks}
The scored synoptic panel consists of the 667 stations for which fixed ERA5
thresholds are available. Country counts are in
Table~\ref{tab:stationcounts}; the full station catalogue shown in
Figure~\ref{fig:network} is larger and is not the scored count.

\begin{table}[h]
\caption{Synoptic stations contributing to the fixed-climatology analysis.}
\label{tab:stationcounts}
\centering\small\begin{tabular}{@{}lr@{}}
\toprule
Country & Stations with fixed thresholds \\
\midrule
AT & 8 \\
BE & 12 \\
CH & 16 \\
CZ & 11 \\
DE & 67 \\
DK & 12 \\
ES & 40 \\
FR & 227 \\
GB & 74 \\
IT & 81 \\
NL & 13 \\
NO & 31 \\
PL & 75 \\
\midrule
Total & 667 \\
\bottomrule
\end{tabular}

\end{table}

\begin{figure}[h]
\centering
\includegraphics[width=.63\textwidth]{fig1_station_network.png}
\caption{Full station catalogues available to the verification system. Only
stations with fixed climatological thresholds enter the reported wind and
temperature regimes; the solar catalogue is likewise larger than its scored
regime panel.}
\label{fig:network}
\end{figure}

Track A covers 1 September 2025--30 June 2026 over DE, FR, GB, ES, IT, PL,
NL, BE, CH, NO, AT, CZ, and DK. It uses all four daily initialization cycles
and the common six-hourly lead grid from 6 to 48\,h. Track B covers
March--June 2026 and compares the regional systems on their common hourly
lead grid.

Solar radiation uses hourly leads from 1 to 48\,h. Regime-stratified results
use 543 stations with fixed thresholds in BE, CH, DE, DK, ES, FR, IT, NL, and
NO. Austrian catalogue stations were not part of the frozen score extraction;
GB and PL are absent, and CZ lacks fixed solar thresholds. Spain contributes
one station through March 2026, and Italian coverage ends after February.

\subsection{Observed-value regimes}
For each station, calendar day, and UTC hour, ERA5 P5, P25, P75, and P95 are
estimated over 1991--2020 with a cyclic $\pm15$-day window. The five regimes
are assigned from the verifying observation. Because they follow the
day-of-year climatology, they represent local anomalies rather than fixed
physical severity.

\section{Complete Track A results}\label{app:tracka}

\begin{table}[h]
\caption{Bias-corrected Track A MAE skill (\%) against IFS, 6--48\,h, by
wind regime. Parentheses give leave-one-month-out jackknife standard errors.}
\label{tab:skillwind}
\centering\small\resizebox{\textwidth}{!}{\begin{tabular}{@{}lrrrrrr@{}}
\toprule
Model & All & $<$P5 & P5--25 & P25--75 & P75--95 & $>$P95 \\
\midrule
Jua EPT-2.1 Europa & +8.4\,(0.5) & -1.7\,(1.4) & +6.4\,(0.9) & +11.3\,(0.7) & +8.5\,(0.4) & +9.0\,(1.2) \\\addlinespace[0.35em]
Jua EPT-2 HRRR & +4.5\,(0.3) & +2.7\,(1.0) & +1.4\,(0.9) & +5.1\,(0.3) & +4.9\,(1.0) & +6.1\,(0.9) \\\addlinespace[0.35em]
Jua EPT-2 Reasoning & +3.0\,(0.1) & -0.7\,(0.4) & +2.8\,(0.5) & +5.6\,(0.3) & +4.5\,(0.3) & -0.7\,(0.3) \\\addlinespace[0.35em]
Jua EPT-2e & -0.0\,(0.2) & -9.1\,(1.0) & -1.4\,(1.1) & +5.6\,(0.5) & +2.9\,(0.6) & -6.6\,(0.5) \\\addlinespace[0.35em]
Microsoft Aurora & -0.7\,(0.3) & -6.6\,(1.1) & -3.9\,(1.1) & +1.9\,(0.8) & +2.6\,(0.8) & -3.4\,(0.6) \\\addlinespace[0.35em]
ECMWF AIFS & -5.4\,(0.4) & -9.8\,(1.2) & -8.0\,(1.3) & -2.1\,(0.6) & -1.4\,(0.8) & -10.6\,(0.6) \\\addlinespace[0.35em]
ECMWF ENS (mean) & +1.4\,(0.2) & -0.9\,(0.7) & +1.5\,(0.6) & +3.1\,(0.4) & +2.1\,(0.4) & -1.0\,(0.2) \\\addlinespace[0.35em]
NOAA GFS & -11.7\,(0.7) & +1.8\,(1.4) & -9.5\,(1.4) & -21.8\,(1.4) & -16.3\,(1.3) & +0.6\,(0.7) \\\addlinespace[0.35em]
DWD ICON Global & +3.6\,(0.4) & -4.1\,(1.2) & +2.8\,(1.0) & +2.6\,(0.7) & +2.5\,(0.7) & +9.5\,(1.1) \\
\bottomrule
\end{tabular}
}
\end{table}

\begin{table}[h]
\caption{Bias-corrected Track A MAE skill (\%) against IFS, 6--48\,h, by
temperature regime.}
\label{tab:skilltemp}
\centering\small\resizebox{\textwidth}{!}{\begin{tabular}{@{}lrrrrrr@{}}
\toprule
Model & All & $<$P5 & P5--25 & P25--75 & P75--95 & $>$P95 \\
\midrule
Jua EPT-2.1 Europa & +11.6\,(0.8) & +5.8\,(1.0) & +10.8\,(1.2) & +10.9\,(0.9) & +11.8\,(1.2) & +15.2\,(1.5) \\\addlinespace[0.35em]
Jua EPT-2 HRRR & +12.2\,(0.4) & +10.6\,(2.4) & +11.5\,(1.1) & +9.5\,(0.7) & +11.4\,(0.8) & +19.7\,(2.3) \\\addlinespace[0.35em]
Jua EPT-2 Reasoning & +4.6\,(0.4) & -5.7\,(1.3) & +2.5\,(1.2) & +7.3\,(0.4) & +7.0\,(0.8) & +1.1\,(1.0) \\\addlinespace[0.35em]
Jua EPT-2e & +0.9\,(0.7) & -21.5\,(3.3) & -5.1\,(2.8) & +5.9\,(0.7) & +7.3\,(1.1) & -5.1\,(0.9) \\\addlinespace[0.35em]
Microsoft Aurora & -1.2\,(1.1) & -20.0\,(4.3) & -7.1\,(4.6) & +3.9\,(1.0) & +3.7\,(1.6) & -6.9\,(1.4) \\\addlinespace[0.35em]
ECMWF AIFS ENS (mean) & +0.4\,(0.5) & -21.3\,(2.6) & -2.8\,(1.9) & +4.5\,(0.8) & +6.1\,(0.9) & -5.1\,(1.7) \\\addlinespace[0.35em]
ECMWF AIFS & -1.3\,(0.6) & -23.9\,(2.7) & -5.0\,(2.0) & +2.0\,(0.8) & +4.6\,(1.5) & -5.0\,(2.3) \\\addlinespace[0.35em]
ECMWF ENS (mean) & -1.7\,(0.3) & -9.8\,(0.9) & -3.3\,(0.8) & +1.3\,(0.3) & +0.2\,(0.8) & -6.1\,(0.8) \\\addlinespace[0.35em]
NOAA GFS & -23.8\,(1.4) & -22.5\,(3.2) & -26.3\,(2.7) & -26.9\,(2.5) & -18.0\,(2.4) & -22.9\,(2.1) \\\addlinespace[0.35em]
DWD ICON Global & +5.6\,(0.9) & -3.6\,(1.3) & +3.3\,(1.6) & +4.1\,(1.3) & +8.4\,(0.7) & +10.2\,(1.9) \\
\bottomrule
\end{tabular}
}
\end{table}

\begin{table}[h]
\caption{Bias-corrected Track A solar MAE skill (\%) against IFS,
1--48\,h, by regime.}
\label{tab:skillsolar}
\centering\small\resizebox{\textwidth}{!}{\begin{tabular}{@{}lrrrrrr@{}}
\toprule
Model & All & $<$P5 & P5--25 & P25--75 & P75--95 & $>$P95 \\
\midrule
Jua EPT-2.1 Helios & +10.2\,(1.7) & +15.4\,(5.0) & +16.4\,(3.4) & -8.2\,(1.1) & +13.7\,(2.9) & +24.8\,(5.4) \\\addlinespace[0.35em]
Jua EPT-2.1 Europa & +7.1\,(1.0) & +3.1\,(4.0) & +11.0\,(3.1) & +9.2\,(2.3) & +1.9\,(4.3) & +4.4\,(4.8) \\\addlinespace[0.35em]
Jua EPT-2 HRRR & +0.6\,(1.4) & -3.6\,(4.6) & +3.8\,(3.7) & +4.5\,(1.0) & -4.6\,(4.8) & -4.1\,(4.9) \\\addlinespace[0.35em]
Jua EPT-2 Reasoning & +6.3\,(0.5) & +0.9\,(1.3) & +9.5\,(0.8) & +14.0\,(0.7) & -0.8\,(1.9) & -2.6\,(2.6) \\\addlinespace[0.35em]
Jua EPT-2e & +2.0\,(1.4) & -10.6\,(3.3) & +5.9\,(2.3) & +13.5\,(1.6) & -7.5\,(4.1) & -7.6\,(6.0) \\\addlinespace[0.35em]
NOAA GFS & -11.2\,(1.6) & -3.3\,(4.4) & -24.3\,(2.1) & -28.3\,(2.5) & +8.8\,(3.4) & +16.7\,(4.2) \\\addlinespace[0.35em]
DWD ICON Global & +5.6\,(0.8) & +6.5\,(3.3) & +8.1\,(2.6) & +3.1\,(1.5) & +4.2\,(3.3) & +6.9\,(4.2) \\
\bottomrule
\end{tabular}
}
\end{table}

\begin{table}[h]
\caption{Track A MAE skill (\%) by lead for 10\,m wind.}
\label{tab:leadwind}
\centering\small\resizebox{\textwidth}{!}{\begin{tabular}{@{}l*{8}{r}@{}}
\toprule
& \multicolumn{2}{c}{6\,h} & \multicolumn{2}{c}{12\,h} & \multicolumn{2}{c}{24\,h} & \multicolumn{2}{c}{48\,h} \\
\cmidrule(lr){2-3} \cmidrule(lr){4-5} \cmidrule(lr){6-7} \cmidrule(lr){8-9}
Model & All & $>$P95 & All & $>$P95 & All & $>$P95 & All & $>$P95 \\
\midrule
EPT-2.1 Europa & +7.6 & +7.8 & +8.0 & +8.4 & +7.7 & +8.2 & +7.3 & +8.0 \\\addlinespace[0.3em]
EPT-2 HRRR & +3.3 & +5.0 & +3.8 & +5.7 & +4.5 & +6.1 & +5.0 & +6.9 \\\addlinespace[0.3em]
EPT-2 Reasoning & +1.9 & -0.8 & +2.3 & -0.9 & +2.9 & -1.0 & +3.9 & -0.6 \\\addlinespace[0.3em]
EPT-2e & -1.3 & -8.4 & -1.1 & -7.9 & -0.3 & -7.1 & +1.0 & -5.9 \\\addlinespace[0.3em]
Aurora & -1.7 & -4.2 & -1.7 & -3.8 & -1.1 & -3.7 & +0.3 & -3.0 \\\addlinespace[0.3em]
AIFS ENS & -3.7 & -8.6 & -3.5 & -9.0 & -2.8 & -9.4 & -0.8 & -8.9 \\\addlinespace[0.3em]
AIFS & -6.8 & -11.5 & -6.9 & -11.8 & -6.0 & -11.9 & -4.2 & -10.3 \\\addlinespace[0.3em]
ENS & +0.2 & -1.1 & +0.5 & -1.1 & +1.1 & -1.1 & +2.3 & -1.7 \\\addlinespace[0.3em]
GFS & -11.2 & +1.2 & -11.3 & +1.0 & -12.0 & -0.1 & -12.4 & -1.2 \\\addlinespace[0.3em]
ICON Global & +4.6 & +9.7 & +4.6 & +9.7 & +4.0 & +9.2 & +2.5 & +8.2 \\
\bottomrule
\end{tabular}
}
\end{table}

\begin{table}[h]
\caption{Track A MAE skill (\%) by lead for 2\,m temperature.}
\label{tab:leadtemp}
\centering\small\resizebox{\textwidth}{!}{\begin{tabular}{@{}l*{8}{r}@{}}
\toprule
& \multicolumn{2}{c}{6\,h} & \multicolumn{2}{c}{12\,h} & \multicolumn{2}{c}{24\,h} & \multicolumn{2}{c}{48\,h} \\
\cmidrule(lr){2-3} \cmidrule(lr){4-5} \cmidrule(lr){6-7} \cmidrule(lr){8-9}
Model & All & $>$P95 & All & $>$P95 & All & $>$P95 & All & $>$P95 \\
\midrule
EPT-2.1 Europa & +8.3 & +12.7 & +12.0 & +16.0 & +11.5 & +14.7 & +11.7 & +14.6 \\\addlinespace[0.3em]
EPT-2 HRRR & +10.1 & +19.6 & +12.6 & +21.4 & +12.1 & +19.5 & +12.3 & +18.5 \\\addlinespace[0.3em]
EPT-2 Reasoning & +2.2 & -0.5 & +3.5 & +0.2 & +4.5 & +0.9 & +6.1 & +2.3 \\\addlinespace[0.3em]
EPT-2e & -4.1 & -10.0 & -1.5 & -7.5 & +1.2 & -4.8 & +3.2 & -2.8 \\\addlinespace[0.3em]
Aurora & -4.5 & -7.7 & -2.2 & -7.0 & -1.1 & -7.0 & +0.3 & -6.7 \\\addlinespace[0.3em]
AIFS ENS & -4.0 & -10.2 & -1.0 & -6.6 & +0.3 & -4.9 & +2.5 & -3.3 \\\addlinespace[0.3em]
AIFS & -4.5 & -7.6 & -2.4 & -5.3 & -1.6 & -5.1 & +0.4 & -3.8 \\\addlinespace[0.3em]
ENS & -2.9 & -6.8 & -2.5 & -6.3 & -1.8 & -6.1 & -0.5 & -5.7 \\\addlinespace[0.3em]
GFS & -27.4 & -24.6 & -24.3 & -22.7 & -24.3 & -23.3 & -22.4 & -22.8 \\\addlinespace[0.3em]
ICON Global & +5.3 & +10.2 & +7.0 & +11.7 & +6.2 & +10.3 & +4.1 & +8.8 \\
\bottomrule
\end{tabular}
}
\end{table}

\begin{table}[h]
\caption{Track A MAE skill (\%) by lead for solar radiation.}
\label{tab:leadsolar}
\centering\small\resizebox{\textwidth}{!}{\begin{tabular}{@{}l*{8}{r}@{}}
\toprule
& \multicolumn{2}{c}{6\,h} & \multicolumn{2}{c}{12\,h} & \multicolumn{2}{c}{24\,h} & \multicolumn{2}{c}{48\,h} \\
\cmidrule(lr){2-3} \cmidrule(lr){4-5} \cmidrule(lr){6-7} \cmidrule(lr){8-9}
Model & All & Ovc. & All & Ovc. & All & Ovc. & All & Ovc. \\
\midrule
EPT-2.1 Helios & +10.8 & +19.1 & +9.0 & +15.2 & +8.1 & +15.4 & +9.1 & +15.5 \\\addlinespace[0.3em]
EPT-2.1 Europa & +5.7 & +10.1 & +6.8 & +9.9 & +6.2 & +10.2 & +6.4 & +10.2 \\\addlinespace[0.3em]
ICON Global & +5.7 & +9.8 & +6.1 & +8.7 & +5.3 & +8.1 & +4.3 & +6.3 \\\addlinespace[0.3em]
EPT-2 Reasoning & +4.2 & +7.0 & +5.1 & +8.5 & +6.0 & +9.5 & +7.8 & +11.1 \\\addlinespace[0.3em]
EPT-2e & -4.5 & +2.4 & -2.5 & +2.1 & +0.1 & +2.9 & +4.1 & +8.1 \\\addlinespace[0.3em]
EPT-2 HRRR & -4.8 & -0.5 & -2.6 & +1.2 & -1.7 & +2.6 & +1.0 & +5.1 \\\addlinespace[0.3em]
GFS & -13.8 & -30.2 & -12.5 & -29.5 & -11.8 & -25.8 & -11.1 & -24.3 \\
\bottomrule
\end{tabular}
}
\end{table}

\section{Robustness checks}\label{app:robustness}

\subsection{Raw and corrected forecasts}
Tables~\ref{tab:robustwind} and \ref{tab:robusttemp} show that the
domain-mean correction materially changes wind-tail rankings and reverses the
sign of the AIFS temperature-tail comparison. The tail-minus-all column is
the direct test used in the main paper.

\begin{table}[h]
\caption{Track A wind MAE skill (\%): raw and bias-corrected all-condition and
$>$P95 values. Tail $-$ all is computed after correction; parentheses give
its monthly jackknife standard error.}
\label{tab:robustwind}
\centering\small\resizebox{\textwidth}{!}{\begin{tabular}{@{}lrrrrr@{}}
\toprule
Model & Raw all & Raw $>$P95 & Corrected all & Corrected $>$P95 & Tail $-$ all \\
\midrule
EPT-2.1 Europa & +0.4 & -5.8 & +7.6 & +8.1 & +0.5\,(0.8) \\
EPT-2 HRRR & +4.1 & +4.0 & +4.4 & +6.2 & +1.8\,(0.7) \\
EPT-2 Reasoning & +2.8 & -0.2 & +2.9 & -0.8 & -3.7\,(0.5) \\
EPT-2e & -1.0 & -7.0 & -0.2 & -6.9 & -6.8\,(0.8) \\
Aurora & -1.0 & -1.9 & -0.9 & -3.5 & -2.7\,(0.9) \\
AIFS ENS & -2.1 & -4.9 & -2.4 & -9.1 & -6.7\,(0.9) \\
AIFS & -6.5 & -12.3 & -5.7 & -11.3 & -5.6\,(1.0) \\
ENS & +3.2 & +4.7 & +1.3 & -1.2 & -2.5\,(0.4) \\
GFS & -8.4 & +13.1 & -11.9 & -0.1 & +11.8\,(1.3) \\
ICON Global & +0.2 & -0.5 & +3.7 & +9.1 & +5.4\,(1.0) \\
\bottomrule
\end{tabular}
}
\end{table}

\begin{table}[h]
\caption{As Table~\ref{tab:robustwind}, for temperature.}
\label{tab:robusttemp}
\centering\small\resizebox{\textwidth}{!}{\begin{tabular}{@{}lrrrrr@{}}
\toprule
Model & Raw all & Raw $>$P95 & Corrected all & Corrected $>$P95 & Tail $-$ all \\
\midrule
EPT-2.1 Europa & +12.3 & +18.1 & +11.6 & +15.2 & +3.6\,(1.2) \\
EPT-2 HRRR & +8.1 & +15.8 & +12.2 & +19.7 & +7.5\,(2.1) \\
EPT-2 Reasoning & +3.9 & -0.2 & +4.6 & +1.1 & -3.5\,(0.9) \\
EPT-2e & +1.1 & -5.4 & +0.9 & -5.1 & -5.9\,(1.2) \\
Aurora & -3.4 & -9.7 & -1.2 & -6.9 & -5.6\,(1.8) \\
AIFS ENS & +4.7 & +2.9 & +0.4 & -5.1 & -5.5\,(1.7) \\
AIFS & +3.4 & +4.4 & -1.3 & -5.0 & -3.7\,(2.0) \\
ENS & -4.4 & -11.4 & -1.7 & -6.1 & -4.4\,(0.6) \\
GFS & -28.1 & -29.0 & -23.8 & -22.9 & +0.9\,(2.3) \\
ICON Global & +9.1 & +16.9 & +5.6 & +10.2 & +4.7\,(1.7) \\
\bottomrule
\end{tabular}
}
\end{table}

\subsection{Threshold definition}
The variants use the same forecast pairs but different regime constructions.
The fixed definition uses station--day--hour anomaly thresholds; the
in-window definition uses one country-level percentile distribution over
stations and times and therefore selects absolute tails. The large
climatological lower-temperature-tail losses of several learned systems
collapse under the in-window definition, and several tail-minus-all signs
also change.

\begin{table}[h]
\caption{Bias-corrected tail skill (\%) under fixed 1991--2020 climatological
thresholds (Clim.) and percentiles of the evaluation window (Window).}
\label{tab:thresholdsensitivity}
\centering\small\resizebox{\textwidth}{!}{\begin{tabular}{@{}lrrrrrr@{}}
\toprule
 & \multicolumn{2}{c}{Wind $>$P95} & \multicolumn{2}{c}{Temperature $>$P95} & \multicolumn{2}{c}{Temperature $<$P5} \\
\cmidrule(lr){2-3}\cmidrule(lr){4-5}\cmidrule(lr){6-7}
Model & Clim. & Window & Clim. & Window & Clim. & Window \\
\midrule
EPT-2.1 Europa & +8.1 & +4.9 & +15.2 & +14.3 & +5.8 & +10.6 \\
EPT-2 HRRR & +6.2 & +2.2 & +19.7 & +12.7 & +10.6 & +11.7 \\
EPT-2 Reasoning & -0.8 & +0.4 & +1.1 & +6.9 & -5.7 & +4.3 \\
EPT-2e & -6.9 & -3.8 & -5.1 & -0.5 & -21.5 & +0.3 \\
Aurora & -3.5 & -1.0 & -6.9 & -0.5 & -20.0 & -1.5 \\
AIFS ENS & -9.1 & -6.9 & -5.1 & -6.3 & -21.3 & -0.3 \\
AIFS & -11.3 & -6.1 & -5.0 & -11.7 & -23.9 & -2.0 \\
ENS & -1.2 & -0.3 & -6.1 & -3.3 & -9.8 & -2.0 \\
GFS & -0.1 & +3.6 & -22.9 & -31.1 & -22.5 & -25.8 \\
ICON Global & +9.1 & +5.9 & +10.2 & +10.1 & -3.6 & +3.9 \\
\bottomrule
\end{tabular}
}
\end{table}

\section{Regional and country-level results}\label{app:regional}

\begin{figure}[h]
\centering
\includegraphics[width=.9\textwidth]{fig8b_track_b_bars.png}
\caption{Track B bias-corrected MAE skill against IFS by
observation-conditioned regime, March--June 2026, 1--48\,h. Uncertainty uses
four monthly blocks.}
\label{fig:trackbbars}
\end{figure}

\begin{table}[h]
\caption{Track B wind MAE skill (\%), 1--48\,h, by regime.}
\label{tab:trackbwind}
\centering\small\resizebox{\textwidth}{!}{\input{tables/t3_track_b_wind_148.tex}}
\end{table}

\begin{table}[h]
\caption{Track B temperature MAE skill (\%), 1--48\,h, by regime.}
\label{tab:trackbtemp}
\centering\small\resizebox{\textwidth}{!}{\begin{tabular}{@{}lrrrrrr@{}}
\toprule
Model & All & $<$P5 & P5--25 & P25--75 & P75--95 & $>$P95 \\
\midrule
Jua EPT-2.1 Europa & +13.4\,(0.6) & +7.9\,(0.9) & +11.9\,(1.5) & +12.5\,(1.3) & +13.5\,(0.4) & +16.8\,(2.0) \\\addlinespace[0.35em]
Jua EPT-2 HRRR & +12.2\,(0.6) & +14.6\,(2.3) & +15.0\,(1.2) & +9.2\,(0.3) & +8.8\,(1.0) & +17.5\,(3.3) \\\addlinespace[0.35em]
DWD ICON-EU & +12.5\,(0.8) & +9.5\,(0.7) & +11.9\,(1.4) & +10.7\,(1.2) & +11.3\,(1.4) & +17.4\,(2.8) \\
\bottomrule
\end{tabular}
}
\end{table}

\begin{table}[h]
\caption{Track B solar MAE skill (\%), 1--48\,h, by regime.}
\label{tab:trackbsolar}
\centering\small\resizebox{\textwidth}{!}{\begin{tabular}{@{}lrrrrrr@{}}
\toprule
Model & All & $<$P5 & P5--25 & P25--75 & P75--95 & $>$P95 \\
\midrule
Jua EPT-2.1 Helios & +12.6\,(2.5) & +21.6\,(5.2) & +18.1\,(5.1) & -11.3\,(0.6) & +13.5\,(5.9) & +32.2\,(4.6) \\\addlinespace[0.35em]
Jua EPT-2.1 Europa & +7.8\,(2.0) & +7.4\,(4.8) & +12.8\,(3.0) & +10.7\,(1.8) & +0.1\,(4.7) & +3.4\,(8.4) \\\addlinespace[0.35em]
Jua EPT-2 HRRR & +0.8\,(2.3) & -0.2\,(5.6) & +4.8\,(5.4) & +6.5\,(0.3) & -6.1\,(9.1) & -6.4\,(8.4) \\\addlinespace[0.35em]
DWD ICON-EU & +7.2\,(1.4) & +12.4\,(4.1) & +12.4\,(2.2) & +1.0\,(2.3) & +0.9\,(3.3) & +10.5\,(6.2) \\
\bottomrule
\end{tabular}
}
\end{table}

\section{Additional methodological details}\label{app:method}

For each model, valid week, initialization hour, and lead, the additive
correction is estimated by pooling station errors from the preceding four
valid weeks. It is one Europe-wide scalar, not a station-specific correction.
Regime labels depend only on the observation and are unchanged by correction.

If $\hat S$ is the full estimate and $\hat S_{(-j)}$ omits month $j$, the
reported standard error is
\[
\operatorname{SE}_{\mathrm{JK}} =
\sqrt{\frac{M-1}{M}\sum_{j=1}^{M}
\left(\hat S_{(-j)}-\overline{\hat S}_{(-\cdot)}\right)^2}.
\]
MAE is pooled linearly by sample count in both the point estimate and each
leave-one-month-out replicate. Track A has $M=10$ and Track B $M=4$.
These values describe sensitivity to months; they are not exact confidence
intervals and do not capture country-level dependence.
